\chapter{Introduction}\label{ch:intro}

This report is a long, multi-file workload for the OverBranch compiler
benchmark. Its purpose is to exercise the parts of a \LaTeX{} build that
only matter at scale: several input files, a table of contents, cross-references
that need more than one pass to settle, and enough pages that the cost of
each run is clearly measurable.

\section{Motivation}\label{sec:motivation}

Long documents are common in practice. Theses, technical reports and
manuals routinely run to dozens of pages, split into chapters that are
maintained by different people. A compiler that handles such documents well
must resolve labels across files, keep the table of contents consistent, and
converge on a stable set of page numbers.

\lipsum[1-2]

\section{Scope}\label{sec:scope}

The report covers a simple statistical model, the estimators used to fit it,
the algorithms that compute those estimators, and a set of synthetic
experiments. Chapter~\ref{ch:background} gives the background, and
Chapter~\ref{ch:model} introduces the model itself. The experiments in
Chapter~\ref{ch:experiments} rely on the algorithms of
Chapter~\ref{ch:algorithms}.

\lipsum[3-4]

\section{Organisation}\label{sec:organisation}

Each chapter is stored in its own file under \texttt{sections/}. The main file
simply inputs them in order. Equation numbers, figure numbers and table
numbers are all automatic, so a reader can follow the references without
manual bookkeeping.

\lipsum[5-6]

\section{Further remarks}\label{sec:further-introduction}

The remarks in this section restate the main point of the chapter in more
general terms, so that the argument can be read without the surrounding
derivations.

\lipsum[51-62]
